\section{An uncut surface estimate from complete segment traces}
\label{sec:v88-surface}
The negative residence moment acts on the original volume source.
To apply it to the outgoing surface one must still pay the change
of the physical angular density. Integrating over the entire segment
provides precisely the factor that makes this payment finite without
an initial-gradient cutoff.

For $1<q<\infty$ define, for a measurable roof interval $J$,
\begin{equation}\label{eq:v88-weak-norm}
 \mathcal N_q(f;J)=\sup_{E\subset J,\ 0<|E|<\infty}
     |E|^{-1+1/q}\int_E|f(t)|\,dt.
\end{equation}
This is a norm equivalent to the usual weak $L^q$ norm. We shall
use its subadditivity directly, not a change of density representative.
Write $J^{[s]}=J+(-s,s)$.

\begin{lemma}[A complete interval-pulse trace]
\label{lem:v88-pulse}
Let a measure space $\Sigma$ carry measurable finite intervals
$(a_\sigma,b_\sigma)$ of lengths $0<\ell_\sigma\le s$ and
nonnegative weights $T_\sigma$. Put
$f(t)=\int_\Sigma T_\sigma\one_{(a_\sigma,b_\sigma)}(t)\,d\sigma$.
If $\upsilon=1-1/q$, then
\begin{equation}\label{eq:v88-pulse}
 \mathcal N_q(f;J)\le
  \int_{\{(a_\sigma,b_\sigma)\cap J\ne\varnothing\}}
                         T_\sigma\ell_\sigma^{1-\upsilon}\,d\sigma.
\end{equation}
The assertion holds also for extended nonnegative integrals.
\end{lemma}
\begin{proof}
For $|E|=v$, Tonelli bounds the contribution of a segment by
$T_\sigma\min(v,\ell_\sigma)$. In both cases $v\le\ell_\sigma$
and $v>\ell_\sigma$,
$v^{-\upsilon}\min(v,\ell_\sigma)\le\ell_\sigma^{1-\upsilon}$.
Integrate this inequality before taking the supremum over $E$.
No disjointness of the roof intervals is assumed or needed.
\end{proof}

\begin{theorem}[Weighted first-exit trace on every short segment]
\label{thm:v88-weighted-trace}
Use the complete Borel segment bijection in the original exact
event and all its nongrazing outgoing endpoints. Let
$T_j(\vartheta)=A_j(b_j-,\vartheta)$, $\ell_j=b_j-a_j$, and
\[
 Z_s(t)=\sum_j\int_{E_j}
          T_j(\vartheta)\one_{\{0<\ell_j\le s,\ a_j<t<b_j\}}\,d\vartheta.
\]
For $0<\eta<1/2$, $\kappa=(1-\eta)/2$ and
$0<\upsilon<\min(\gamma_0,1-\kappa)$, set $q=(1-\upsilon)^{-1}$.
At each finite count,
\begin{equation}\label{eq:v88-surface-before-limit}
 \mathcal N_q(Z_s;J)
 \le\int_{J^{[s]}}q_{\upsilon,s}(u)\,du
       +C_\eta s^{1-\kappa-\upsilon}
                              \int_{J^{[s]}}v^\eta(u)\,du.
\end{equation}
Thus for each fixed $L,s>0$,
\begin{align}
 \limsup_m\sup_{R,n,k}\sup_{|J|=L}
                      \mathcal N_q(m^2Z_s;J)
 \le{}& C_\upsilon(L+2s)s^{\gamma_0-\upsilon}\notag\\
      &+C_\eta(L+2s)s^{1-\kappa-\upsilon}.
                                      \label{eq:v88-short-weak}
\end{align}
All positive residence lengths, all initial gradients and all incoming
endpoint types are included; no optical-sensitivity cutoff is used.
\end{theorem}
\begin{proof}
Fix one selected segment and suppress its angular index. Its exact
physical density satisfies $|A'|\le C A I_m/z$. For $a<t<b$,
\[
 A(b-)\le A(t)+C\int_t^b A(u)I_m(u)/z(u)\,du.
\]
Integrate over the entire interval $(a,b)$, multiply by
$\ell^{-\upsilon}$, and use Tonelli. This gives
\begin{equation}\label{eq:v88-trace-triangle}
 A(b-)\ell^{1-\upsilon}\le
  \int_a^b A(t)\ell^{-\upsilon}\,dt
   +C\ell^{-\upsilon}\int_a^b (u-a)A(u)I_m(u)/z(u)\,du.
\end{equation}
Forward radial growth, first from an interior point and then in the
limit toward $a$, implies $z(u)^2\ge2\lambda_\circ(u-a)$. Hence
\[
 (u-a)z(u)^{-1}
 \le C_\eta(u-a)^{1-\kappa}z(u)^{-\eta}
 \le C_\eta\ell^{1-\kappa}z(u)^{-\eta}.
\]
The second term of \eqref{eq:v88-trace-triangle} is therefore at
most $C_\eta s^{1-\kappa-\upsilon}\int_a^b A I_mz^{-\eta}$.
This calculation pays the variation of both $\omega$ and $\chi$
in $A=\omega/(4\pi c_*\chi)$. It does not replace an inverse
Jacobian in a first-exit formula by one.

Every selected segment that intersects $J$ lies in $J^{[s]}$.
Sum \eqref{eq:v88-trace-triangle} over precisely these physical
segments. The Borel bijection and its incident-component convention
count each interior point once. The first sum is bounded by the
actual source $q_{\upsilon,s}$; the second by the actual fractional
source $v^\eta$. Nongrazing outgoing traces are used on the left;
dropping that endpoint selector only on the two positive upper
sources introduces neither a new normalization nor another label.
Lemma~\ref{lem:v88-pulse} proves \eqref{eq:v88-surface-before-limit}.

Initially all statements may be read as nonnegative extended
inequalities on compact subsegments. The finite actual moments,
Tonelli and positive exhaustion make the right side finite and
justify the complete sum and endpoint limit. Apply
\eqref{eq:v88-negative-ordered} and
\eqref{eq:v83-fixed-width-input}, fixing both interval lengths and
all bands before the collision limsup, to obtain
\eqref{eq:v88-short-weak}.
\end{proof}

\begin{theorem}[The complete untruncated first-exit surface]
\label{thm:v88-full-surface}
For every $1<q<32/31$, put $\upsilon=1-1/q$ and
$\nu_q=\gamma_0-\upsilon=1/q-31/32>0$. For every fixed $L>0$,
\begin{equation}\label{eq:v88-full-surface}
 \limsup_m\sup_{R,n,k}\sup_{|J|=L}
       \mathcal N_q(m^2\overline{\mathcal F}_r;J)
                                  \le C_{q,L}r^{\nu_q}
\end{equation}
for all sufficiently small previously fixed $r>0$. This is the
actual one-sided outgoing first-exit surface with no lower residence
or gradient cutoff. For example $q=64/63$ gives $\nu_q=1/64$.
\end{theorem}
\begin{proof}
The complete residence kernel splits into its long part $\mathcal L_r$
with $\ell>2r$ and a short part bounded by $Z_{2r}$. The retained
long-residence height theorem \eqref{eq:v88-long-input} gives
$\mathcal N_q(m^2\mathcal L_r;J)\le |J|^{1/q}m^2\|\mathcal L_r\|_\infty$
and an ordered upper bound $CL^{1/q}r^{\gamma_0}$.
Use Theorem~\ref{thm:v88-weighted-trace} with $s=2r$ for the
complete short part. Take, for instance, $\eta=1/4$. Then
$1-\kappa-\upsilon>\gamma_0-\upsilon>0$.
Subadditivity of \eqref{eq:v88-weak-norm} proves the result.
The short contribution was summed at finite count through an actual
negative moment, not reconstructed from infinitely many band limsups.
\end{proof}

The weak exponent is finite. Equation~\eqref{eq:v88-full-surface}
controls the distribution of heights and does not assert that the
surface essential supremum tends to zero. The positive cutoff has
been removed from this new norm theorem, not from the earlier
essential-height theorem.
